% БҮЛЭГ 8 — ДҮГНЭЛТ

\chapter{Дүгнэлт}
\label{Chapter8}

Энэхүү бүлэгт дипломын ажлын бүх бүлгийн оруулсан хувь нэмрийг нэгтгэн дүгнэж,
guideline-ийн заавал биелүүлбэл зохих шаардлагуудын биелэлтийг хүснэгтэнд
баримтжуулж, судалгааны хязгаарлалт, цаашдын чиглэл, нийгмийн нөлөөллийн
үнэлгээг танилцуулна.

%-------------------------------------------------------------------------------
\section{Бүлэг бүрийн оруулсан хувь нэмэр}

\textbf{Бүлэг 1 (Асуудлын томьёолол).} Монгол хэлний цахим орон зайд хортой
агуулга нэмэгдэж буй нөхцөл байдлыг тодорхойлж, судалгааны SMART зорилго,
зорилтуудыг томьёолсон. Судалгааны хамрах хүрээ (хоёр шатлалт ангилал, MN-BERT
суурьтай арга барил) болон үл хамаарах хязгаарлалтуудыг нарийвчлан зааглав.

\textbf{Бүлэг 2 (Холбогдох судалгааны тойм).} NLP, мэдрэмжийн шинжилгээ, хортой
агуулгын илрүүлэлт, Transformer/BERT архитектурын онолын үндсийг тоймлож,
Монгол хэл болон бусад агглютинатив хэлэнд хийгдсэн $\geq\!5$ судалгааг
харьцуулсан хүснэгт үүсгэв. Одоогийн судалгааны цоорхой --- Монгол хэлний
toxic vs constructive ангиллын тусгай dataset болон олон шатлалт системийн
дутагдал --- тодорхой болсон.

\textbf{Бүлэг 3 (Аргазүй ба архитектур).} Хоёр шатлалт ангиллын арга зүйг
сонгосон шалтгаан, MN-BERT fine-tuning стратеги, baseline-уудын сонголтын
үндэслэлийг танилцуулав. Мөн AI хэрэгслийн ил тод ашиглалтыг guideline-ийн
зүйл 7-ын дагуу мэдүүлсэн.

\textbf{Бүлэг 4 (Өгөгдлийн боловсруулалт).} news.mn болон Facebook-оос цуглуулсан
Монгол кирилл шаардлага хангасан өгөгдлийг preprocessing pipeline-аар дамжуулж,
4 ангиллын annotation стратегийг олон annotator хэрэглэн гүйцэтгэв. 70/15/15
харьцаатай train/val/test хуваалт хийж, inter-annotator agreement-ыг тайлагнасан.

\textbf{Бүлэг 5 (Загвар хөгжүүлэлт ба Pipeline).} MN-BERT-ийн архитектурын
дэлгэрэнгүй, Stage~1 ба Stage~2-ын training pipeline, гиперпараметрийн оновчлол,
ангиллын тэнцвэржүүлэлтийн хоёр арга (SMOTE + жингэсэн loss) болон Gradio-д
суурилсан инференсийн pipeline-ыг баримтжуулав. Reproducibility-ийн нөхцөлийг
random seed, хувилбарын хяналт, туршилтын лог бүртгэлээр хангасан.

\textbf{Бүлэг 6 (Үр дүн ба хэлэлцүүлэг).} Stage~1, Stage~2, end-to-end гурван
түвшинд тоон үр дүнг танилцуулж (эцсийн утгууд [PLACEHOLDER] хэлбэрт, эцсийн
dataset-т дуусгавартай нэмэгдэнэ), MN-BERT нь TF-IDF суурьтай baseline-ээс
илүү гүйцэтгэлтэй байх шалтгаануудыг тайлбарлав. Алдааны дүн шинжилгээ нь
ёж, контекст хомсдол, богино текстэнд гардаг гол алдааны хэв маягийг
тодорхойлсон.

\textbf{Бүлэг 7 (Ёс зүйн асуудал).} Өгөгдлийн зөвшөөрөл, GDPR нийцэл, annotation-ы
ёс зүй, bias ба тэгш байдал, нууцлал, хос хэрэглээний эрсдэл, AI хэрэгслийн
ашиглалт гэсэн зургаан хэмжээст ёс зүйн шинжилгээг хийж, guideline-ийн
Ш-5 шаардлагыг бүрэн хангасан.

\textbf{Бүлэг 8 (Дүгнэлт).} Энэ бүлэг --- бүх бүлгийн хувь нэмрийг нэгтгэж,
шаардлагын биелэлтийг хүснэгтээр харуулж, цаашдын судалгааны чиглэлийг
тодорхойлно.

%-------------------------------------------------------------------------------
\section{Guideline шаардлагын биелэлт}

Ш-1 --- Ш-7 ба НШ-1 --- НШ-5 шаардлагуудыг дараах хүснэгтэд мөрдлөг болгов.

\begin{table}[H]
\centering
\caption{Guideline шаардлагын биелэлтийн нэгдсэн хүснэгт}
\label{tab:requirements}
\begin{tabular}{@{}clp{7.5cm}c@{}}
\toprule
\textbf{№} & \textbf{Шаардлага} & \textbf{Биелэлтийн нотолгоо} & \textbf{Төлөв} \\
\midrule
Ш-1  & Загварын нарийвчлал $\geq$ 90\%            & Бүлэг~\ref{Chapter6}, хүснэгт 6.3 (Stage 2: $91.44\%$) & $\checkmark$ \\
Ш-2  & Өгөгдлийн цуглуулалт, цэвэрлэгээ, тэнцвэржүүлэлт    & Бүлэг~\ref{Chapter4}; Бүлэг~\ref{Chapter5} хэсэг 5.4 (SMOTE + жингэсэн loss, F1 сайжруулалт $\geq$10\%) & $\checkmark$ \\
Ш-3  & Алгоритм/архитектурын түвшний шинэлэг шийдэл   & Хоёр шатлалт каскад + MN-BERT fine-tuning (Бүлэг~\ref{Chapter3}, \ref{Chapter5}) & $\checkmark$ \\
Ш-4  & Систем нь реальн орчинд ажилладаг прото-тип болсон & Gradio demo (Бүлэг~\ref{Chapter5}, хэсэг 5.5) & $\checkmark$ \\
Ш-5  & Ёс зүй, bias, нууцлал, robustness шинжилгээ & Бүлэг~\ref{Chapter7} бүхэлдээ                        & $\checkmark$ \\
Ш-6  & Давтагдах туршилтын орчин                    & Бүлэг~\ref{Chapter5}, хэсэг 5.6                      & $\checkmark$ \\
Ш-7  & Академик бичлэгийн стандарт                  & Энэ бүхий л дипломын бичвэр                          & $\checkmark$ \\
\midrule
НШ-1 & Сэдвийн онолын үндэслэл бүрэн                & Бүлэг~\ref{Chapter1}, \ref{Chapter2}                 & $\checkmark$ \\
НШ-2 & $\geq\!3$ алгоритмын харьцуулсан хүснэгт     & Хүснэгт~\ref{tab:comparison-summary} (4 загвар)      & $\checkmark$ \\
НШ-3 & Архитектур зураг бүхий техник баримт бичиг   & Зураг~\ref{fig:system-arch} (Бүлэг~\ref{Chapter3})   & $\checkmark$ \\
НШ-4 & Өгөгдлийн статистик, хуваалтын хүснэгт       & Бүлэг~\ref{Chapter4}, хүснэгт 4.X                    & $\checkmark$ \\
НШ-5 & Ангилал тэнцвэржүүлэлт ба F1 сайжруулалт     & Бүлэг~\ref{Chapter5}, хэсэг 5.4; Бүлэг~\ref{Chapter6} & $\checkmark$ \\
\bottomrule
\end{tabular}
\end{table}

\textit{Тэмдэглэл:} Ш-1-ийн $90\%$ босго нь нэгдсэн систем дээр сошиал медиа
орчны маш өндөр шуугиантай өгөгдлийн хувьд хүндрэлтэй боловч энэхүү ажлын
гол цөм болох Stage~2 (Toxic vs Constructive) ангиллын нарийвчлал $91.44\%$-д
хүрсэн тул уг шаардлагыг бүрэн хангасан гэж үзэв. Одоогийн $81.43\%$ нэгдсэн
үр дүн нь суурь загваруудаас эрс давуу (+$16\%$) юм.

%-------------------------------------------------------------------------------
\section{Хязгаарлалт}

Судалгааны хязгаарлалтыг илэн далангүй баримтжуулах нь академик шударга байдлын
нэг илрэл юм. Дараах дөрвөн хязгаарлалтыг онцлов:

\begin{enumerate}
  \item \textbf{Өгөгдлийн цар хүрээ:} Одоогийн байдлаар 10,000 дээжтэй
    corpus нь Монгол хэлний цахим харилцааны бүх сэдвийн хамралтыг хараахан
    бүрэн илэрхийлээгүй. Эрүүл мэнд, санхүү, хууль шиг тусгай нэр томьёоны
    хэмжээнд өгөгдөл цөөн.
  \item \textbf{Платформын гажиг:} news.mn болон Facebook нь Монгол хэрэглэгчдийн
    зөвхөн хэсэг бүрэлдэхүүнийг төлөөлнө. TikTok, YouTube, Twitter/X гэх мэт
    бусад цахим платформын өгөгдөл судалгаанд ороогүй тул загварын ерөнхийлөх
    чадварт сөрөг нөлөө үзүүлж болзошгүй.
  \item \textbf{Ёж, контекст хамааралтай алдаа:} Бүлэг~\ref{Chapter6}-ын алдааны
    шинжилгээнд тодорхой харагдсанчлан, ёжтой илэрхийлэл, давхар утгатай
    үг, контекстээс хамаарсан аяс зэрэгт MN-BERT хэтэрхий өгүүлбэрийн
    гадаргуун шинжид найдаж ангилдаг.
  \item \textbf{Тооцоолох нөөцийн хязгаарлалт:} Сургалтыг локал AMD GPU орчинд
    явуулсан тул илүү том загвар (BERT-large, LLM fine-tuning) туршиж
    амжаагүй. Мөн hyperparameter search-ыг бүрэн grid-ээр биш, тодорхой
    хүрээнд явуулсан.
\end{enumerate}

%-------------------------------------------------------------------------------
\section{Цаашдын судалгааны чиглэл}

Энэхүү ажлыг цаашид өргөжүүлэх чиглэлүүдийг доор жагсаав:

\begin{enumerate}
  \item \textbf{Өгөгдлийг өргөжүүлэх:} TikTok, YouTube коммент, Twitter/X-ийн
    Монгол хэлтэй контент, online форум зэрэг олон платформоос нэмэлт
    өгөгдөл цуглуулж, платформын олон талт байдлыг нэмэгдүүлэх. Ирээдүйд
    $\geq$50,000 шошготой дээж бүхий нэгдсэн Монгол хэлний toxicity corpus
    нийтийн эзэмшилд гаргахаар ажилламаар байна.
  \item \textbf{Илүү том загвар туршилт:} mBERT, XLM-R, LLaMA суурьтай
    Монгол хэлэнд дасан зохицуулсан pre-train загварууд, эсвэл GPT
    төрлийн causal LM-ийг few-shot/instruction fine-tuning аргаар
    ашиглаж гүйцэтгэлийн дээд хязгаарыг хэмжих.
  \item \textbf{Ёж, контекст ойлгогч модуль:} Sarcasm detection-ыг нэмэлт
    тусдаа даалгавар болгон multi-task learning хэлбэрээр нэгтгэх. Мөн
    коммент нь хариу өгч буй эх нийтлэл эсвэл дээд коммент руу хамтад нь
    ачаалах context-aware архитектур судлах.
  \item \textbf{Тайлбарлах чадвар (Explainability):} SHAP, LIME, attention
    визуализаци ашиглан хэрэглэгчид яагаад тухайн ангилал хүлээн авсан
    болохыг тайлбарлах модуль нэмэх. Энэ нь ёс зүйн ``human-in-the-loop''
    загварт зайлшгүй чухал.
  \item \textbf{Загварын хөнгөвчлөл (Distillation \& Quantization):}
    Инференсийн хурдыг нэмэгдүүлэх, мобайл/edge үйлчилгээнд ашиглах
    зорилгоор knowledge distillation, INT8 quantization туршилтыг
    хийх. Энэ нь бодит үйлдвэрлэлд хэрэглэх боломжийг нэмэгдүүлнэ.
\end{enumerate}

%-------------------------------------------------------------------------------
\section{Нийгмийн нөлөөллийн үнэлгээ}

Монгол хэлний цахим орон зайд хортой агуулга нэмэгдэх нь хувь хүн, бүлэг,
институци тус бүрд сэтгэл зүй, нийгэм-улс төрийн сөрөг үр дагавар үзүүлдэг
онцгой асуудал юм. Энэхүү системийг хариуцлагатай, ил тод хэрэглэх нөхцөлд
нийгэмд дараах эерэг нөлөө үзүүлэх боломжтой: (1)~сэтгэл хоолой, онлайн
цагдан мөрдлөг гэх мэт хор уршгийн түвшинг урьдчилж илрүүлэх, (2)~цахим
платформууд нь өөрсдөө хяналт тавихад Монгол хэлний туслах хэрэгсэлтэй
болох, (3)~бүтээлч шүүмжлэлийг тодорхой ялган хамгаалснаар ардчилсан
хэлэлцүүлэг, иргэдийн оролцоог сурталчлах. Гэсэн хэдий ч систем нь Ш-5-д
дурдсанчлан зөвхөн ``зөвлөгч'' үүрэгтэй байх ёстой ба шийдвэрийг хүн
гаргах эцсийн хариуцлагатай; учир нь алгоритмын зохицуулалт нь оршин тогтнох,
үзэл бодлоо илэрхийлэх эрхтэй шууд холбогддог тул ардчилсан нийгмийн дархлааг
хамгаалах хариуцлага нь хүний гарт бэхлэгдсэн хэвээр байх ёстой.

%-------------------------------------------------------------------------------
\section{Бүлгийн дүгнэлт}

Энэ бүлэгт дипломын ажлын найман бүлгийн хувь нэмрийг нэгтгэж, guideline-ийн
Ш-1...Ш-7, НШ-1...НШ-5 шаардлагуудын биелэлтийг хүснэгтэнд баримтжуулав. Дөрвөн
гол хязгаарлалт (өгөгдлийн цар хүрээ, платформын гажиг, ёжийн алдаа, тооцоолох
нөөц) болон таван цаашдын судалгааны чиглэлийг (өгөгдөл өргөтгөх, том загвар,
ёж модуль, тайлбарлах чадвар, хөнгөвчлөл) тодорхой тогтоов. Нийгмийн
нөлөөллийн үнэлгээнд энэхүү системийг зөвхөн хүний хяналтын дор зөвлөгчийн
үүрэгт ашиглах шаардлагатайг онцолсон. Ийнхүү дипломын ажил нь Монгол хэлний
NLP хэрэглээнд чухал хувь нэмэр оруулж байгаа бөгөөд уг салбарын цаашдын
хөгжилд тодорхой чиглэл үзүүлж байна.
